\section{Response geometry, rank-one costs, and remaining boundaries}\label{s6:synthesis}

The preceding algorithms avoid enumerating an unrestricted parametric LP
response. This distinction is necessary even on a physical blending path.
We first construct such paths and a one-pool extension, then examine the
cost structure of an isolated rank-one block. The examples separate three
questions: the size of a complete response description, exact representation
of a nonconvex set, and the complexity of obtaining one optimizer.

\subsection{A telescoping certificate on a path}\label{s6:telescoping}
For $0<\epsilon<1/2$, let
\begin{equation}
 P_\epsilon=\{x\in\R^n:0\le x_1\le1,\quad
 \epsilon x_{j-1}\le x_j\le1-\epsilon x_{j-1}\ (2\le j\le n)\},
 \qquad x_0=0.\label{s6:km}
\end{equation}
Every coordinate lies in $[0,1]$. The two inequalities in each pair
cannot both be tight. Thus a vertex has one tight inequality in each
pair; conversely any such choice gives a nonsingular triangular system
and a feasible vertex. Its coordinates are
$x_j(u)=u_j+(1-2u_j)\epsilon x_{j-1}(u)$ for $u\in\{0,1\}^n$.
All $2^n$ terminal values are distinct: the last bit chooses the disjoint
intervals $[0,\epsilon]$ and $[1-\epsilon,1]$, and within either interval
its affine recursion can be inverted, successively recovering all bits.

\begin{lemma}[A vertex certificate]\label{s6:certificate}
Put
\[
 d_j=(1-\epsilon)\epsilon^{2(n-j)-1}\quad(j<n),\qquad
 L(x)=x_n-\sum_{j<n}d_jx_j,\qquad F(x)=L(x)-x_n^2.
\]
Then
\begin{equation}
 F(x)=\sum_{j=1}^n\epsilon^{2(n-j)}
 (x_j-\epsilon x_{j-1})(1-\epsilon x_{j-1}-x_j)\ge0
 \quad(x\in P_\epsilon),\label{s6:telescoping-identity}
\end{equation}
with equality precisely at the vertices. Every vertex $v$ is uniquely
exposed by $\sum_{j<n}d_jx_j+(2v_n-1)x_n$.
\end{lemma}
\begin{proof}
The $j$th product expands to
$x_j-x_j^2-\epsilon x_{j-1}+\epsilon^2x_{j-1}^2$.
The weighted squares telescope, leaving $-x_n^2$, and the remaining
coefficients give $L$. All factors and weights are nonnegative, so equality
requires a tight endpoint in every pair, exactly the vertex condition.
For $t=v_n$, throughout $P_\epsilon$,
\[
 2tx_n-L(x)\le 2tx_n-x_n^2=t^2-(x_n-t)^2\le t^2.
\]
Equality requires $F=0$ and $x_n=t$, which select $v$ uniquely.
\end{proof}

Exponential two-dimensional shadows of this Klee--Minty cube are established
by G\"artner, Helbling, Ota and Takahashi
\cite[Section~4]{s6:gartner2013}. The lemma supplies a convenient
self-contained certificate for the physical constructions below; it is
not a claim that the exponential-shadow phenomenon is new.
For comparison, their direction at $\epsilon=1/4$ uses
$c_j=\epsilon^{3(n-j)}$ for $j<n$, $c_n=0$, and exposes every vertex
by $c^Tx+\lambda x_n$ with $|\lambda|<1/15$.
The direction $d$ in the lemma instead uses the interval $[-1,1]$.
Both directions have polynomial rational encoding length.

\subsection{Physical paths and one varying product price}\label{s6:physical-paths}
Fix $n\ge2$, $\epsilon=1/4$, and $s_j=4^{j-n}$. Source $j$ has
quality $C_j=n-j$ and exact supply $s_j$. Its only arcs go to products
$j-1$ and $j$, with flows $s_j-t_j$ and $t_j$, each bounded above by
$s_j$. Product~0 has capacity $s_1$ and redundant upper quality $C_1$;
product~$n$ has capacity one and redundant upper quality $C_n=0$.
For $2\le j\le n$, product $j-1$ has capacity $s_j$ and upper quality
$(C_{j-1}+C_j)/2$.

\begin{proposition}[Exact path realization]\label{s6:path-realization}
This pool-free physical network is affinely bijective to $P_{1/4}$ by
$t_j=s_jx_j$. Its bipartite bypass graph is one path, and every source
and product has degree at most two. All upper flow bounds are at most
one; scaling all qualities by $1/(n-1)$ puts them in $[0,1]$.
\end{proposition}
\begin{proof}
The internal product receives $t_{j-1}+s_j-t_j$. Its capacity is
$t_j\ge t_{j-1}$. Subtracting its midpoint quality from the two
incident source qualities gives opposite coefficients $1/2$ and
$-1/2$, so its homogeneous quality row is
$t_{j-1}\le s_j-t_j$. Thus its two rows are exactly
\begin{equation}
 t_{j-1}\le t_j\le s_j-t_{j-1}.\label{s6:physical-strip}
\end{equation}
After dividing by $s_j$, the predecessor coefficient is $1/4$.
The endpoint and arc bounds give the other bounds, and the reverse
substitution satisfies every physical row. All other flows are uniquely
$s_j-t_j$; hence the correspondence is bijective. Scaling quality does
not change any homogeneous quality inequality.
\end{proof}

\begin{theorem}[Exponential physical price response]\label{s6:price-response}
There is a family of rational blending paths with $O(n)$ nodes, one
upper-bounded quality, all upper flow capacities at most one, zero
production costs, and positive product revenues, for which varying just
the final product's unit revenue exposes $2^n$ distinct unique optimal
flows. The family can have only two distinct input qualities and no
positive lower flow bounds. Its optimal-value function has $2^n$
open affine pieces. Every fixed-price problem is an LP.
\end{theorem}
\begin{proof}
First keep the exact source contracts in \cref{s6:path-realization}.
Use the direction $d$ from \cref{s6:certificate}; its coefficient in
physical rightward flows is $w_j=d_j/s_j=3s_j$ for $j<n$.
Put $w_n=\lambda\in[-1,1]$ and define output revenues
\[
 R_0=1,\qquad R_j=1+\sum_{h=1}^j w_h\quad(1\le j\le n).
\]
Since $\sum_{j<n}3s_j=1-s_1<1$, all revenues lie in $[0,3]$ and only
$R_n$ depends on $\lambda$. With zero costs their total revenue is
\begin{equation}
 \sum_j[R_{j-1}(s_j-t_j)+R_jt_j]
 =C_0+\sum_{j<n}d_jx_j+\lambda x_n,
 \quad C_0=\sum_jR_{j-1}s_j.\label{s6:price-identity}
\end{equation}
The constant $C_0$ is independent of $\lambda$. The lemma exposes every
vertex, including $x_n=0,1$ at $\lambda=-1,1$. Strict optimality among a
finite vertex set persists on a neighborhood of each exposing parameter,
so each has a nonempty one-sided interval inside $[-1,1]$, and an open
interval in its interior. Distinct terminal values give distinct slopes;
there are exactly $2^n$ affine regimes. One can instead use the published
direction $c$ and $w_j=c_j/s_j=16^{j-n}$, obtaining revenues in $[0,2]$
and the smaller price interval $[-1/15,1/15]$.

To use only two input qualities, replace the descending quality path by
the following relay path. Start with a quality-one source of supply
$s_1$. For each $j=2,\ldots,n$, attach a quality-zero source of supply
$s_j$ through a product with capacity $s_j$ and upper quality $1/2$.
If $j<n$, append a quality-one source of the same supply $s_j$ through
a product with exact demand $s_j$ and redundant upper quality one.
A descending step gives \eqref{s6:physical-strip}. At a reset product,
\[
 t_j+(s_j-\widetilde t_j)=s_j
 \quad\Longleftrightarrow\quad\widetilde t_j=t_j.
\]
Thus the quality-one duplicate copies the signal before the next
quality-zero source. End products have redundant upper quality one
and capacities $s_1,1$. There are $m=2n-2$ sources with supply sequence
\[
 S_1,\ldots,S_m=(s_1,s_2,s_2,s_3,s_3,\ldots,s_{n-1},s_{n-1},s_n).
\]
All arcs of source $i$ have upper bound $S_i$.
The original free rightward flows and their equal duplicates determine
all physical flows, giving the same affine bijection to $P_{1/4}$.
Assign coefficient $3s_j$ to the free-coordinate position $j<n$,
zero to reset duplicates, and $\lambda$ to the last position.
With these coefficients in path order denoted $g_i$, take
$\rho_0=1$, $\rho_i=1+\sum_{h\le i}g_h$.
The same telescoping revenue computation gives
$\sum_i\rho_{i-1}S_i+\sum_{j<n}d_jx_j+\lambda x_n$.
Only the final revenue varies and every base revenue is in $[0,3]$.

It remains to remove all source and reset-product lower contracts.
Let $Q$ be the nonempty exact-contract polytope in its $N=2m$ arc
variables $f$. Write these contracts $Ef=e$. Retain their upper bounds
and drop their lower bounds, so the deficit vector
$\Delta=e-Ef$ is nonnegative; put $\delta=\mathbf1^T\Delta$.
All coefficient rows, after clearing the midpoint factor two, have
entries in $\{-1,0,1\}$. By \cref{s3:hoffman}, every relaxed flow has
an exact-contract repair $f'\in Q$ with
$\|f'-f\|_1\le H\delta$, where $H=N^N$.
The base revenue vector has infinity norm at most three, uniformly in
$\lambda$. Choose $M=3H+1$, add $M$ to every output revenue, and add
another $M$ to every reset-output revenue. Conservation makes the first
bonus reward all source throughputs; the second rewards reset throughputs.
Writing $E_0=\mathbf1^Te$, the new objective $B(f)+M(E_0-\delta)$ obeys
\[
 B(f)+M(E_0-\delta)
 \le B(f')+ME_0-(M-3H)\delta.
\]
Thus every optimizer satisfies every removed contract. The value function
changes by exactly $ME_0$ and every previously exposed optimizer remains
unique. Revenues are now positive, all production costs remain zero,
and $M$ has $O(n\log n)$ bits. No bound on revenue magnitude is claimed.
For the original direction $c$, the same proof uses $M=2H+1$.
The version with varying input qualities needs no reset contracts and
uses $N=2n$ and the identical source-throughput bonus argument.
\end{proof}

\begin{corollary}[Explicit elimination needs exponential output]\label{s6:response-size}
For either response in \cref{s6:price-response}, a piecewise-affine table
requires $2^n$ pieces. Any quantifier-free Boolean formula in price and
value alone describing its graph, or its epigraph on the price interval,
uses polynomials whose total degree is at least $2^n$.
The fixed-objective state message
$G(t)=\max\{\sum_{j<n}d_jx_j:x\in P_{1/4},\ x_n=t\}$ has
$2^n-1$ affine intervals.
\end{corollary}
\begin{proof}
Remove identically zero polynomials from a proposed formula. Where none
of its remaining polynomials vanishes, all signs, and hence its truth
value, are locally constant. Each open affine graph segment is a
boundary segment, so the product of the defining polynomials vanishes
on that segment. Restriction to its supporting line vanishes on an
interval and therefore identically; the line equation divides the
product. The $2^n$ distinct lines require that many linear factors.
This applies also after any upper clipping of the epigraph that leaves
all those segments below the clip.
For $G$, all projected vertex points $(x_n,\sum d_jx_j)$ are uniquely
exposed with coefficient $+1$ on the second coordinate. They lie on its
upper boundary, have distinct first coordinates, and consecutive points
cannot be collinear. Their first coordinates run from zero to one,
giving exactly $2^n-1$ intervals.
\end{proof}

These are limits on explicit descriptions. Backward elimination still
evaluates the support at a rational price with $O(n)$ arithmetic operations:
set $a_n=\lambda$ and $a_j=d_j-\epsilon|a_{j+1}|$ for $j=n-1,\ldots,1$;
the value is $\sum_j\max(0,a_j)$. Eliminating a variable with positive
current coefficient selects its upper endpoint, adds that coefficient
to the constant, and subtracts $\epsilon$ times its absolute value from
the predecessor coefficient; a negative coefficient selects its lower
endpoint and makes the same predecessor update without the constant.
All intermediate bit lengths are polynomial. The same recurrence works
with $c$ instead of $d$. Compact circuits, extended formulations, and
implicit algorithms are therefore fully compatible with the lower bound.
In particular, \cref{s4:path-projection} projects onto endpoint coordinates,
whereas $G$ retains a dense eliminated objective.

\begin{proposition}[A limit on individual coordinate ports]\label{s6:ports}
Without nonphysical global rows, a pool-free network with source and
product degrees at most two cannot represent every bounded rational
polytope by a coordinate projection onto selected individual arc flows,
even with unrestricted attribute count and two-sided bounds.
\end{proposition}
\begin{proof}
Every physical row involves at most two arc variables. Fourier--Motzkin
elimination preserves this property: a pair of rows of opposite signs
in the eliminated variable contains at most one other variable each.
Thus every coordinate projection has a finite two-variable-per-inequality
description; its size need not be polynomial. The simplex
$\{x\in\R_+^3:x_1+x_2+x_3\le1\}$ has no such description.
The point $(1/2,1/2,1/2)$ satisfies every valid inequality involving at
most two coordinates, because each of its coordinate pairs extends to
a simplex point by setting the third coordinate to zero. Yet that point
is outside the simplex. This does not restrict arbitrary linear images.
\end{proof}

\subsection{One pool: many optima and an exact LP hull}\label{s6:geometry}
Use the contracted path of \cref{s6:path-realization}, retaining its
unscaled qualities $C_j=n-j$, and put $t=t_n=x_n$. Modify its terminal
product and add the following interface:
\[
 A\to n,\quad A\to P,\quad B\to P,\quad B\to W,\quad
 P\to W,\quad P\to V,\quad Z\to V.
\]
Sources $A,B$ have qualities one and two and exact supplies one;
source $Z$ has quality zero and supply interval $[0,1]$. Every new arc
has capacity one. Pool $P$ has exact throughput one. Products $n,W$
have exact demand one and upper qualities one and two; product $V$ has
capacity two and upper quality one. There are no further positive
lower bounds. The original terminal source has quality zero.

\begin{theorem}[A physical vertex-forcing family]\label{s6:vertex-forcing}
The resulting one-pool network has a feasible set affinely equivalent,
in physical arc flows and active pool quality, to
\begin{equation}
 T=\{(x,z):x\in P_{1/4},\quad x_n-x_n^2\le z\le1\}.
 \label{s6:T}
\end{equation}
There are ordinary nonnegative source costs and product revenues bounded
by two for which the network has exactly $2^n$ isolated global maximizers.
Every source and product has degree at most two, the pool has two feeds
and two outlets, and the full undirected network is connected with one
cycle. All quality data can be scaled into $[0,1]$.
\end{theorem}
\begin{proof}
Demand at $n$ forces $f_{An}=1-t$; source $A$ then gives $f_{AP}=t$.
Pool throughput gives $f_{BP}=1-t$, whence source $B$ gives $f_{BW}=t$.
Demand at $W$ forces $f_{PW}=1-t$, leaving $f_{PV}=t$.
The pool is active with quality $q=2-t$. If $z=f_{ZV}$, the only
nonredundant new quality row is
\[
 (2-t)t\le t+z\quad\Longleftrightarrow\quad z\ge t-t^2.
\]
Every $t\in[0,1]$ has $0\le t-t^2\le1/4$, so all other arc,
source, throughput and quality bounds hold, and this describes every
physical extension. In particular all flows and $q$ are affine in
$(x,z)$, despite the curved feasible boundary.

Give path source $j$ cost $s_j$, sources $A,B$ cost one, and $Z$ cost
two. Give path product $j<n$ revenue $s_{j+1}$ and products $n,W,V$
revenue one. Subtracting one from every source cost and every product
revenue leaves profit unchanged by total mass conservation. The new
interface prices are zero except for unit cost on $Z$. On the path,
source $j$ has cost $s_j-1$, its left output has the same revenue,
and its right output has revenue $s_{j+1}-1$ when $j<n$ and zero
when $j=n$. Thus profit is exactly
\begin{equation}
 \sum_{j<n}(s_{j+1}-s_j)t_j-z
 =\sum_{j<n}3s_j^2x_j-z=\sum_{j<n}d_jx_j-z.
 \label{s6:geometry-profit}
\end{equation}
By \cref{s6:certificate}, this is at most $-F(x)\le0$; equality forces
a path vertex and the unique clean flow $z=t-t^2$. These are all
$2^n$ optimizers, so all are isolated. The degree assertions follow from
the displayed arcs. The original path and its attached trees introduce
no cycle; the sole cycle is $B-P-W-B$. Dividing every quality and
specification by $\max(n-1,2)$ preserves all flows and normalizes data.
\end{proof}

\begin{corollary}[Strict local optima with distinct values]\label{s6:local}
A polynomially encoded nonnegative price perturbation of this network
has at least $2^n$ strict local maxima with distinct values, exactly one
of the constructed points being globally optimal. Its prices are at
most three.
\end{corollary}
\begin{proof}
Increase both the revenue of $V$ and cost of $Z$ by $\delta=4^{-n}$.
Their net contribution is $\delta[(t+z)-z]=\delta t$.
At least clean flow, the objective is $H(x)=-F(x)+\delta x_n$.
At a vertex of terminal value $t$ it is $\delta t$. Every terminal
vertex value has denominator dividing $4^{n-1}$, and these values are
distinct. For neighboring vertices $v,v'$ with terminal values $t,t'$,
linearity of $L$ and $F(v)=F(v')=0$ give
\[
 F((1-a)v+av')=a(1-a)(t'-t)^2.
\]
The outward derivative of $H$ is
$-(t'-t)^2+\delta(t'-t)<0$, since $|t'-t|\ge4^{-(n-1)}>\delta$.
The tangent cone at a polytope vertex is generated by its incident edge
rays. Strict negativity on these finitely many rays bounds the linear
term above by $-\eta\|x-v\|$ for some $\eta>0$ on that cone; the
quadratic remainder is $O(\|x-v\|^2)$. Hence the vertex is a strict
local maximum. Extra clean flow decreases profit strictly, so the
property transfers to \eqref{s6:T} and then to the physical coordinates.
Finally every feasible profit is at most $\delta t\le\delta$, with
equality only at the unique vertex with $t=1$ and least clean flow.
No assertion excludes additional local maxima. The perturbation shrinks
with dimension and supplies no uniform tolerance or local-solver runtime
bound.
\end{proof}

\begin{corollary}[Exact integer dimension of the optimal set]\label{s6:integer-lift}
Let $S$ be the complete optimal physical-flow set for the unperturbed
objective in \cref{s6:vertex-forcing}, optionally retaining the pool quality.
The minimum number of integer coordinates in any exact convex lift of
$S$ is $n$, allowing arbitrary continuous auxiliaries and unbounded
integer variables.
\end{corollary}
\begin{proof}
For distinct optimal path vertices $v,v'$,
$F((v+v')/2)=(v_n-v_n')^2/4>0$, so their physical midpoint is not
optimal. Suppose a convex set with $k<n$ integer coordinates projects
exactly to $S$ when those coordinates are integral. Choose a lift of
each of its $2^n$ points. Two integer vectors have the same parity;
their midpoint is integral and belongs to the convex lifting set.
Its physical projection is the forbidden midpoint, a contradiction.
For the upper bound, attach to the optimizer indexed by $u$ the vector
$u\in\{0,1\}^n$ and take the convex hull of these finitely many lifted
points. Any integral vector in its final $n$ coordinates is binary;
a convex combination yielding it must use only points with that same
binary vector. It therefore selects exactly its original optimizer.
This upper construction may have exponentially many vertices.
The lower-bound argument is the established midpoint/parity principle
of Lubin, Vielma and Zadik \cite[Section~4.2]{s6:lubin2022}, applied to
this physical optimal set.
\end{proof}

\begin{theorem}[A compact exact hull and polynomial optimization]\label{s6:lp-hull}
The same physical feasible family admits an exact polynomial-size LP
convex hull. Every rational linear objective in physical flows and active
pool quality has a polynomial-time exact physical optimizer, with rational
coordinates of polynomial bit length.
\end{theorem}
\begin{proof}
Writing $d^Tx=\sum_{j<n}d_jx_j$, the exact hull in \eqref{s6:T} is
\begin{equation}
 \conv(T)=\{(x,z):x\in P_{1/4},\quad d^Tx\le z\le1\}.
 \label{s6:hull}
\end{equation}
The certificate gives $d^Tx\le x_n-x_n^2$, proving containment.
Conversely write $x$ as a convex combination of path vertices.
At every vertex $v$, $v_n-v_n^2=d^Tv$, so $(x,d^Tx)$ is the same
convex combination of points in $T$. Also $(x,1)\in T$. Their vertical
segment contains every point on the right of \eqref{s6:hull}.
A vertex of that polytope must lie on $z=d^Tx$ or $z=1$, since otherwise
it admits a vertical perturbation. On either affine boundary, a nonvertex
$x$ admits a nontrivial convex decomposition. Consequently every hull
vertex has $x$ a path vertex and is physically feasible. The two
boundaries do not meet, since $d^Tx\le x_n-x_n^2\le1/4$.
Solve the rational LP and return an optimal vertex, then apply the affine
physical map. An arbitrary nonvertex point of its optimal face need not
be physically feasible. Polynomial bit length follows from rational LP
vertex bounds. The hull does not exactly represent the original
nonconvex set or its finite optimal set, so it does not contradict
\cref{s6:integer-lift}.
\end{proof}

\begin{corollary}[Three qualities and general supply sequences]\label{s6:alphabet-geometry}
The conclusions of \cref{s6:vertex-forcing,s6:local,s6:integer-lift,s6:lp-hull}
hold with only three distinct input qualities. The unperturbed geometry,
integer dimension and LP hull also extend to any rational supplies with
$s_n=1$, $s_1>0$ and $s_{j+1}>2s_j$ in the descending-quality path.
\end{corollary}
\begin{proof}
For the first assertion attach the identical interface to the contracted
two-quality relay path in the proof of \cref{s6:price-response}. Its
last source has quality zero, so every interface calculation is unchanged.
Keep all its source, reset-output and interface contracts. At physical
source position $i$ give cost $S_i$, left endpoint output revenue $S_1$,
and output after position $i<m$ revenue $S_{i+1}$; terminal and interface
prices remain as before. The coefficient of its rightward flow is
$S_{i+1}-S_i$. This coefficient vanishes when the next source is the
same-supply reset duplicate. The coefficient $3s_j$ before the next
descending step occurs on that duplicate, whose rightward flow equals
$t_j$. At the first position it occurs directly on $t_1$. Thus the
sum is again \eqref{s6:geometry-profit}. All conclusions transfer through
the affine bijection, and the full quality alphabet is $\{0,1,2\}$.

For the second assertion keep \eqref{s6:physical-strip}, which remains
the physical capacity/quality pair for these supplies. Direct expansion gives
\[
 t_n-t_n^2-\sum_{j<n}(s_{j+1}-s_j)t_j
 =\sum_{j=1}^n(t_j-t_{j-1})(s_j-t_{j-1}-t_j),\qquad t_0=0.
\]
Every summand is nonnegative. The strict growth of supplies makes both
endpoint branches distinct and their terminal ranges disjoint, so there
are exactly $2^n$ endpoint vertices and distinct terminal values.
The interface and economics now use coefficients $s_{j+1}-s_j$.
The equality, midpoint and convex-hull proofs apply verbatim in these
physical coordinates. Encoding and algorithmic bounds include the given
supply bit lengths. The specific perturbation $4^{-n}$ is asserted only
for the geometric sequence.
\end{proof}

The linear exact-penalty argument in \cref{s6:price-response} does not
remove the contracts in the nonlinear interface. In fact deleting just
the lower pool-throughput bound changes its optimum: set
$x_1=\cdots=x_{n-2}=0$, $x_{n-1}=1$, $x_n=1/4$, send $1/4$ from
$A$ through the pool to $V$, send the rest of $A$ to $n$, send all of
$B$ to $W$, and set $z=f_{PW}=0$. All other bounds and contracts hold,
the pool quality is one, and profit is $d_{n-1}=3/16>0$.
Multiple local optima in pooling have earlier geometric examples;
Grothey and McKinnon \cite[Section~3]{s6:grothey2020} describe a
three-pool, two-quality example. The present family adds the explicit
asymptotic count, degree restrictions and LP hull; this comparison does
not establish exhaustive priority for those refinements.

\subsection{Why one unrestricted aggregate changes the abstract problem}
\label{s6:slab}
The preceding geometry alone cannot prove degree-two pooling hardness.
The following abstract reduction identifies an additional ingredient.
Rank-one concave quadratic hardness in general is established by
Pardalos and Vavasis \cite{s6:pardalos1991}; the point here is its local
path structure and the role of a dense slab.

\begin{proposition}[A path with one dense slab]\label{s6:slab-hardness}
For rational $0<\epsilon<1/2$, deciding whether
\[
 \min\{F(x):x\in P_\epsilon,\ B-1/4\le w^Tx\le B+1/4\}\le0
\]
is ordinarily $\NP$-complete, even on nonempty compact instances.
The Hessian of $F$ has one negative eigenvalue and rank one. Local
path coefficients can alternatively be fixed to $0,\pm1,\pm1/4$,
with additional singleton upper bounds and polynomially many variables.
\end{proposition}
\begin{proof}
For a SUBSET SUM instance of positive integer weights $w_i$ and integer
$0\le B\le W=\sum_iw_i$, set $\epsilon=1/(8W)$. A vertex with
endpoint bits $u$ satisfies $|x_i-u_i|\le\epsilon$.
If $w^Tu=B$, then $|w^Tx-B|\le1/8$, so that vertex satisfies the
slab and $F=0$. Conversely $F\le0$ forces a vertex by
\cref{s6:certificate}; its bits satisfy
$|w^Tu-B|\le1/4+1/8<1$, forcing equality of the integers.
The zero vertex has weighted sum zero. The all-one-bit vertex has
weighted sum at least $W-1/8$. If $B<W$, convex interpolation reaches
$B$; if $B=W$, the latter vertex already lies in the slab.
Thus the domain is nonempty. The coefficients in $F$, including powers
of $\epsilon$, have polynomial bit length. The endpoint bits are an
$\NP$ certificate: reconstruct the rational vertex and verify the slab.
The Hessian is $-2e_ne_n^T$.

To fix $\epsilon=1/4$, let $r$ be the least nonnegative integer with
$4^{-(r+1)}\le1/(8W)$. Place free coordinates at positions
$1,r+2,2r+3,\ldots$ on a path of dimension $N=n+(n-1)r$.
Impose $x_j\le1/2$ on each padding coordinate and use the full
$N$-coordinate certificate $F$. At a zero of $F$, padding positions
must take their lower branch: the upper branch is at least $3/4$.
Successive free coordinates are separated by $r$ multiplications by
$1/4$, so each differs from its endpoint bit by at most $4^{-(r+1)}$.
Every free-bit vector extends uniquely to a zero of $F$. Put the slab
weights only on free positions; the same rounding, nonemptiness and
certificate arguments apply. Since $r=O(\log W)$ this is polynomial.
The dense weights and objective coefficients retain binary numerical
data, so this gives no strong-hardness or inverse-polynomial gap claim.
\end{proof}

Fixing the scalar nonlinear coordinate $q=x_n$ in this problem leaves
a bounded LP. The local inequalities form a path, but the slab has two
bounds on a dense aggregate, and $L$ is another dense linear form.
Thus one nonlinear parameter and a fixed number of aggregate rows do
not justify replacing the independent blocks of \cref{s4:fixed-core}
by arbitrarily long coupled paths. No physical degree-two realization
of the extra weighted slab has been supplied. In its absence, even
minimizing $a^Tx-bx_n^2$ over the bare path, for $b\ge0$, is polynomial:
a concave function has a vertex minimizer, and $x_n^2=L(x)$ at every
vertex, so minimize the linear function $a^Tx-bL(x)$ instead. The slab
creates new vertices and invalidates that linearization.

\subsection{Exact optimization of low-interaction-rank pool costs}
\label{s6:low-rank}
Consider the isolated rank-one margin block
\begin{equation}
 \mathcal K=\{W\in\R_+^{m\times n}:\rank W\le1,\quad
 \ell\le W\mathbf1\le u,\quad
 \ell'\le W^T\mathbf1\le u'\},\label{s6:margin-block}
\end{equation}
with finite nonnegative rational bounds in the stated order, and objective
$\langle C,W\rangle$. For $W\ne0$, its margins $x,y$ and common total
$S$ give $W=xy^T/S$ by \eqref{s1:rank-one-form}.
The feasible-total interval is
$I=[\max(\mathbf1^T\ell,\mathbf1^T\ell'),
\min(\mathbf1^Tu,\mathbf1^Tu')]$; reject an empty interval.
Define the \emph{interaction rank}
\[
 k=\rank D,\qquad D_{ij}=C_{ij}-C_{i1}-C_{1j}+C_{11}.
\]
This is the least rank after subtracting arbitrary additive row and
column costs. Indeed multiplying $C-a\mathbf1^T-\mathbf1b^T$ on the
left by $I_m-\mathbf1e_1^T$ and on the right by $I_n-e_1\mathbf1^T$
gives $D$, so no subtraction can have smaller rank. Taking
$a_i=C_{i1}$, $b_j=C_{1j}-C_{11}$ attains it.
Rational elimination therefore constructs
$C=a\mathbf1^T+\mathbf1b^T+UV^T$ with $k$ factor columns.
The objective is
\begin{equation}
 a^Tx+b^Ty+\frac{(U^Tx)^T(V^Ty)}S.\label{s6:interaction-objective}
\end{equation}

\begin{theorem}[Fixed interaction rank]\label{s6:rank-theorem}
For each fixed $k$, exact linear optimization over \eqref{s6:margin-block}
is polynomial in the rational input bit length, with dimension dependence
$(m+n)^{O(k+1)}$. An optimizer uses rational numbers and at most one
square root, all with polynomial encoding length. This is polynomial
for fixed $k$; the displayed dependence is not an FPT bound in $k$.
\end{theorem}
\begin{proof}
A linear image of an $h$-coordinate box in dimension $d$ is a translated
sum of $h$ segments. Its vertices correspond to full-dimensional sign
regions of the central hyperplanes orthogonal to its nonzero generators.
The region recurrence obtained by adding one hyperplane gives at most
$2\sum_{i=0}^{d-1}\binom{h-1}{i}$ vertices for positive effective
dimension, and one in dimension zero. Degenerate or repeated generators
only decrease this number. For fixed $d$, enumerate these vertices
constructively by adding segments one at a time: union the existing
vertices with their translates and retain only vertices of the resulting
convex hull. Membership of one candidate in the convex hull of the others
is a rational LP test. Store with every retained candidate a corner of
the original box mapping to it. Intermediate lists have the same
polynomial bound, and every stored coordinate is a rational generator
sum of polynomial bit length.

Apply this construction to
\[
 Z_x=\{(\mathbf1^Tx,U^Tx,a^Tx):\ell\le x\le u\},\qquad
 Z_y=\{(\mathbf1^Ty,V^Ty,b^Ty):\ell'\le y\le u'\}.
\]
Both projected boxes have dimension at most $k+2$ and polynomially many
vertices for fixed $k$. At any fixed $S>0$, the objective
\eqref{s6:interaction-objective} is affine in either projected margin
when the other is fixed. Starting with an optimum, minimize first on
one slice at total $S$ and choose a slice vertex, then do the same on
the other. Both steps preserve optimality. There is therefore an optimum
at a pair of slice vertices.

Every vertex of a hyperplane slice of a polytope lies on an original
vertex or edge. To see this, take its smallest containing face. The
point is in the relative interior of that face. If the face dimension
were at least two, intersecting with one hyperplane through the point
would leave a nonzero two-sided feasible direction. This contradicts
extremality in the slice.
Enumerate all pairs of projected-box vertices, including every edge.
A segment with different endpoint totals parametrizes a projected margin
affine in $S$ on the closed interval between those totals. Interpolating
its two stored box corners gives a feasible original-margin witness.
The additional nonedge segments are harmless, since they stay in the
projected box. Include each original vertex as a singleton candidate
at its fixed total. Equal-total segments can be omitted: an interior
point of an edge parallel to the slicing hyperplane cannot be a slice
vertex, and its endpoints are already included.

Pair row and column candidates and intersect their intervals with $I$.
On every positive interval, their projected coordinates and witnesses
are affine in $S$, and substitution gives
\begin{equation}
 f(S)=\alpha S+\beta+\gamma/S\label{s6:fractional-piece}
\end{equation}
with rational coefficients of polynomial bit length. Evaluate positive
endpoints and the stationary point $\sqrt{\gamma/\alpha}$ if
$\alpha,\gamma>0$ and it lies in the interval. These are all interior
minimum possibilities, apart from a constant function, where any
feasible total suffices. When both coefficients are negative a stationary
point is a maximum; the other sign cases are monotone. Singletons are
evaluated directly.

If $0\in I$, zero is feasible and is included separately; if $I=\{0\}$
it is the sole matrix. No missing infimum at zero occurs, since
$|\langle C,W\rangle|\le\|C\|_\infty S$ for every nonnegative
matrix of total $S$. In particular an interval candidate reaching zero
has the correct continuous zero extension. Every enumerated candidate
is feasible, and every fixed-total optimum has representatives on the
list, so minimizing the list gives a global optimum.
There are polynomially many candidates for fixed $k$. Interior values
are $\beta+2\sqrt{\alpha\gamma}$; comparisons of two such values
require sign checks and at most two squarings, or degree-four algebraic
comparison. Interpolation at the chosen total and $W=xy^T/S$ stay in
one quadratic number field. These operations and all rational LP tests
have polynomial bit complexity, proving the stated bound and output
representation.
\end{proof}

Low-rank box optimization and projected-box methods have established
precedents. Punnen, Sripratak and Karapetyan
\cite[Sections~3.2--3.3]{s6:punnen2015} treat bipartite box objectives;
Hlad\'ik, \v{C}ern\'y and Rada \cite{s6:hladik2021} use zonotope faces
for continuous quadratic box optimization. The additional feature here
is the shared variable total and its division in
\eqref{s6:interaction-objective}; the proof handles it by slice segments
and one-variable rational optimization.

\begin{proposition}[Rank-one cost specialization]\label{s6:rank-one-cost}
If $C=pq^T$ is supplied in factored form, the optimum and its margins
can be found in $O((m+n)\log(m+n))$ arithmetic operations and comparisons,
with polynomial bit complexity. Factoring a supplied dense rank-one
matrix or writing every entry of $W$ adds $O(mn)$ arithmetic work.
\end{proposition}
\begin{proof}
At each positive total define $A_-(S),A_+(S)$ as the minimum and maximum
of $p^Tx$ on $[\ell,u]\cap\{\mathbf1^Tx=S\}$, and define $B_-,B_+$
similarly. The two scalar images are intervals, so their product is
minimized at one of the four endpoint pairs, regardless of signs.
The optimal value is
$\min_{\sigma,\tau\in\{-,+\}}A_\sigma(S)B_\tau(S)/S$.
Each endpoint function is continuous knapsack: start at the lower bounds
and fill residual capacities in increasing coefficient order for a
minimum, decreasing order for a maximum. Sorting and prefix sums give
all affine pieces and $O(m+n)$ total breakpoints. Merge those breakpoints
with $I$ and sweep the four paired functions. Each is of form
\eqref{s6:fractional-piece}; process its endpoints and possible minimum
stationary point as above. Recover the selected greedy margins in linear
work. The same exact comparisons prove polynomial bit complexity.
\end{proof}
For example, fix the first row and column margins to one, bound the
other two margins by $[0,1]$, and take $p=(3,1)$, $q=(2,1)$.
The cost is $S+3+2/S$ for $1\le S\le2$, minimized at $S=\sqrt2$.
Even a rank-one cost need not have a rational optimizer.

\begin{corollary}[A fixed-attribute Lagrangian pool oracle]\label{s6:oracle}
Suppose an isolated pool block has additive operating costs, row and
column margin bounds as in \eqref{s6:margin-block}, and $K$ attributes.
For any supplied rational nonnegative multipliers on terminal quality
inequalities, the Lagrangian subproblem obtained by dualizing those
inequalities is exactly solvable in polynomial time for fixed $K$.
\end{corollary}
\begin{proof}
Writing input attributes $Q_{ik}$ and output upper bounds $H_{jk}$,
the homogeneous quality rows are
$\sum_i(Q_{ik}-H_{jk})W_{ij}\le0$.
For minimization, multipliers $\lambda_{jk}\ge0$ add matrix cost
\[
 Q\Lambda^T-\mathbf1 h^T,\qquad h_j=\sum_k\lambda_{jk}H_{jk}.
\]
The additive part does not affect interaction rank, and
$\rank(Q\Lambda^T)\le K$. Apply \cref{s6:rank-theorem}.
Lower quality rows reverse the corresponding signs; both upper and lower
multipliers can be combined into the same $K$ input-attribute columns.
An affine basis of the input qualities can further replace $K$ by its
rank. All remaining constraints must have precisely the margin/rank form;
other pool couplings need separate treatment. This proves an exact
subproblem oracle, with no zero-duality-gap or full pooling tractability
claim.
\end{proof}
If a supplied cost approximation $\widetilde C$ has fixed interaction
rank and $\|C-\widetilde C\|_\infty\le\eta$, where this is the maximum
absolute entry norm, its optimizer has true cost at most
$2\eta S_{\max}$ above optimum, with
$S_{\max}=\min(\mathbf1^Tu,\mathbf1^Tu')$. Indeed every feasible
matrix has objective perturbation at most $\eta S_{\max}$. The bound
assumes both the approximation and its error certificate are supplied.

\subsection{Matrix parameters and nonlinear leaves}\label{s6:parametric}
Interaction rank of a \emph{cost} must be distinguished from rank of a
parameter perturbing a \emph{constraint matrix}.
Consider rational data and
\[
 \min c^Tx\quad\text{subject to }Ax+\lambda Dx\le b,
 \quad x\in X,\quad\underline\lambda\le\lambda\le\overline\lambda,
\]
where $X$ is a rational polyhedron and the nonempty parameter interval
is finite. If $D=uv^T$, introduce $z=v^Tx$ and $w=\lambda z$.
The exact projected feasible set is the union of the two polyhedra
\begin{align*}
 Ax+uw&\le b,\quad x\in X,\quad z=v^Tx,\quad
 z\ge0,\quad\underline\lambda z\le w\le\overline\lambda z;\\
 Ax+uw&\le b,\quad x\in X,\quad z=v^Tx,\quad
 z\le0,\quad\overline\lambda z\le w\le\underline\lambda z.
\end{align*}
For $z\ne0$, recover $\lambda=w/z$; at $z=0$ both branches force
$w=0$ and any allowed parameter works. Thus two LPs solve joint
feasibility and minimization. If the objective is finite and the set
nonempty, a rational LP solution and that ratio give a rational optimum
of polynomial bit length; unbounded objectives are detected by the LPs.
This auxiliary-variable extension is closely related to the
one-nonzero-column case of Boveroux, Carvalho, Lodi and Louveaux's
preprint \cite[Section~3.3]{s6:boveroux2026}.

Their positive-product construction already uses perturbation rank two:
a fixed-one variable $s$ occurs in the opposing rows for
$\lambda s=z_1$, and a second variable $z_2$ in
$\lambda z_2\le t$. The two nonzero columns have independent supports.
Using the positive-product source in \cref{s3:matsui}, with its explicit
bounds, these rows give ordinary hardness for joint minimization of $t$.
This observation concerns dense base-polytope constraints, not physical
path coefficients. A parameter-dependent right-hand side adds a fixed-one
column before measuring perturbation rank; a direct parameter objective
term and a true max--min objective are outside the two-LP argument.

There is also an arithmetic boundary to replacing the polyhedral blocks
in \cref{s4:fixed-core} by convex nonlinear leaves. Given positive
integers $a_i$ and an integer $B$, let
\[
 P_i=\{y_i:0\le y_i\le\max(1,a_i),\quad y_i^2=a_i\}.
\]
Each leaf is a compact convex singleton in one scalar variable, described by
rational quadratic data, and there is no nonlinear core. The single
aggregate inequality $\sum_i y_i\le B$ asks exactly whether
$\sum_i\sqrt{a_i}\le B$. Hence an unrestricted exact polynomial-bit
extension to these leaves would solve Square-Root Sum in polynomial time.
This is an arithmetic implication, not an $\NP$-hardness theorem.
Independent algebraic leaf values may generate a growing common field,
whereas the polyhedral recovery in \cref{s4:fixed-core} takes place in
the field of the fixed core.

\subsection{Related convexification and network results}\label{s6:related}
Rank-one flow formulations and their relaxations are established tools
for pooling \cite{dey2020-convexifications-of-rank-one-based,gupte2017-relaxations-and-discretizations-for-the}.
The unpublished notes cited below are included in this repository. Their
bibliographic paths are relative to the repository root;
\path{papers/pooling/source-index.md} maps citation keys and titles to
source files and their SHA256 hashes.
For the general-cost margin block \eqref{s6:margin-block}, a separate
analysis \cite{s6:margin-note} proves strong $\NP$-completeness of
rational threshold decision, polynomial-size rational threshold witnesses,
quadratic-field optima, and an exact algorithm with dependence
$d2^d\poly(m,n,B)$ for $d=\min(m,n)$ and input bit length $B$.
The hardness reduction has nonadditive matrix costs. Standard source
and product economics gives $C_{ij}=a_i+b_j$, and the quality-free
margin problem is then an LP. Thus this general matrix theorem concerns
optimization and separation of a pool block, and does not by itself
prove physical pooling hardness. The distinction also explains how the
fixed-interaction-rank oracle can coexist with general-cost hardness.

The accompanying exact convexification result \cite{s6:face-note}
uses the unit-margin hull
\[
 K_{2m}=\conv\{W\ge0:\rank W\le1,\ W\mathbf1\le\mathbf1,
 W^T\mathbf1\le\mathbf1\}.
\]
Its affine section
\[
 \operatorname{tr}W=1,\qquad\sum_{i=1}^mW_{i,m+i}=0,
 \qquad\mathbf1^TW\mathbf1=m
\]
is a face affinely isomorphic to
$\operatorname{COR}(m)=\conv\{xx^T:x\in\{0,1\}^m\}$.
The face atoms are
$m^{-1}(x,\mathbf1-x)(x,\mathbf1-x)^T$; the inverse projection is
$X=mW_{[m],[m]}$. At a rank-one generator the trace is at most one,
and trace equality forces both margins to be the same nonempty binary
vector. The zero cross-pair sum excludes choosing both members of a
pair. Maximizing total mass on this face then selects one member of
every pair, proving the assertion. Consequently established correlation
polytope lower bounds transfer by affine section and projection:
Fawzi and Parrilo \cite[Theorem~1]{s6:fawzi2013} treat fixed-order PSD
blocks and SOC conversion, while Lee, Raghavendra and Steurer
\cite[Theorem~1.1]{s6:lrs-full} treat unrestricted PSD order.
They give exponential total second-order-cone dimension bounds,
exponential numbers of fixed-order PSD blocks, and superpolynomial
unrestricted PSD order. Scalar inequalities count as nonnegative cone
factors, or PSD blocks of order one; SDP size is the sum of block orders.
The exact LP lower bound is only a formal consequence: this hull is
already nonpolyhedral for two rows and columns, so it has no finite LP
lift at all. Positive-error LP bounds below carry the useful size claim. They permit arbitrary real coefficients and do not depend
on efficiently constructing the lift.

Approximate formulations require a separate robustness argument.
The stability analysis \cite{s6:stability-note} and its SDP extension
\cite{s6:approx-sdp-note} use entrywise matrix $\ell_1$ error, with
$U_1=\{E:\sum_{i,j}|E_{ij}|\le1\}$:
\[
 K_{2m}\subseteq Q\subseteq K_{2m}+\eta U_1,
 \qquad A_m=m(184m+6).
\]
For $\eta\le1/A_m$, every LP lift of such an outer polyhedron has
$2^{\Omega(m)}$ inequalities. For $\eta\le1/(2A_m)$, every exact
PSD lift of such a convex outer set has total matrix order at least
$\exp(\Omega((m/\log m)^{2/13}))$ for sufficiently large $m$.
The LP transfer uses Braun, Fiorini, Pokutta and Steurer's fixed-factor
sandwich bound \cite[Theorem~6]{s6:bfps2015}. The SDP transfer instead
uses shifted quadratic slacks and the pseudo-density machinery of
Lee, Raghavendra and Steurer
\cite[Theorem~3.8, equation~(3.11), and Theorem~5.3]{s6:lrs-full}. These are uniform outer-approximation
statements at a specified inverse-quadratic accuracy, equivalently
uniform support gaps for all costs with maximum absolute entry at most
one. They do not constrain an objective-specific relaxation, a fixed
cost-rank family, or a weaker accuracy requirement. The conic arguments
are developed in the cited unpublished research notes; their role here
is to delimit exact and uniform approximate pool-block convexification.

A different adjacent program studies common-factor products $w_j=xy_j$.
With finite boxes and a fixed number of linking rows
$Ay+Bw\le c+dx$, linear optimization is polynomial for a scalar $x$,
including signed intervals and an integer scalar
\cite{s6:common-factor-note}. The mechanism is bounded-variable LP
basis enumeration with a fixed number of rows, univariate sign
partitions and derivative-root tests. This is an optimization oracle,
not a general compact conic hull theorem. A reciprocal anchor gives
another useful distinction: the hull of
$(X,1/X,Y,XY)$ for $a\le X\le b$, $0<a<b$, $Y\in[0,1]^n$
requires a common distribution of $X$ across all leaves.
Intersecting individual one-leaf hulls can fail to enforce that
compatibility \cite{s6:anchor-faces-note}. The full hull has a
polynomial rational separation procedure based on a call-function
envelope of affine moment bounds \cite{s6:anchor-full-note}; for consecutive integer
$X$ values, mean-preserving rounding and telescoping sums avoid
enumerating the binary-encoded scalar range
\cite{s6:integer-anchor-note}. These statements impose no additional
product bounds or linking rows on the anchored original points;
intersecting such restrictions with the hull need not remain exact.

Sparse network--simplex products form another model:
$x$ belongs to a bounded equality-constrained network-flow polytope,
$y\ge0$, $\mathbf1^Ty\le1$, and selected products $z_{ej}=x_ey_j$
are retained. Simplex disaggregation already yields a compact extended
hull in general \cite[Appendix, equation~(25)]{s6:khademnia2024v2}.
The locator refers to arXiv:2302.14151v2 (26 February 2024).
The separate universality result
\cite{s6:network-universality-note} concerns coordinate sections and
unbounded original-space facet coefficients at simplex dimension two;
it transfers the established three-way transportation universality of
De Loera and Onn \cite{s6:deloera2006}, and is not an optimization
hardness claim. Cycle and theta blocks instead admit explicit hull
relations through one- and two-dimensional circulation coordinates
\cite{s6:network-cycle-note}. More generally, blocks consisting of
internally disjoint parallel paths admit exact subset cuts with
coefficients in $\{-1,0,1\}$ on flow and product coordinates and
polynomial separation \cite{s6:network-parallel-note}. Simplex
coefficients retain the network data. This graph class can have large
cycle rank within a block, but it is not the class of all series--parallel
graphs; converting inequality balances to slack arcs can change the
required structure. None of these network--simplex statements transfers
to physical pooling without a correspondence preserving mixing and
capacity constraints.

Finally, addition and reciprocal gadgets have a separate application
to resistive power-flow equations \cite{s6:power-note}. The corresponding
AC transfer depends on positive voltage magnitudes, zero reactive
injections, and bounded differences of real phase angles; principal-angle
constraints on cycles have different semantics. These are separate
physical equations and complexity questions. The common reduction
method does not make power-flow claims part of a pooling theorem.

\subsection{What the parameter boundaries leave open}\label{s6:open}
The main remaining structural question is standard pooling with a fixed
number of pools, fixed affine input-quality rank, degree-two bypass
graphs, and unboundedly many pool attachments. General flow and quality
intervals are allowed; feasibility and dense-profit optimization should
be treated separately. Fixing the pool qualities makes each local
external-node fiber small, but its projected polygon depends on the
unknown parameters. The one-parameter bound in \cref{s5:quasipoly-path}
is quasipolynomial in general, not a polynomial symbolic composition
result. Even an efficient endpoint relation would forget the pool-flow
contributions needed in the dense global pool mass and quality balances.
A fixed count of those aggregates alone is insufficient, as
\cref{s6:slab-hardness} shows for an abstract system.

A shared fixed-quality helper pool does not automatically convert the
degree-three copy network into a degree-two bypass network. Each original
midpoint source has its own private supply equation; a shared helper
preserves only their total. For example, in two full midpoint-copy
pairs with endpoint qualities zero and two, midpoint quality one,
and demand four at each copy output, choose both endpoint-pair signals
zero in the first pair and two in the second. The required four midpoint
flows are $4,4,0,0$, totaling eight. A common helper with throughput eight
can supply them, but the original private midpoint supplies would demand
four in each pair. The local replacement therefore loses a required
relation, regardless of whether additional constraints elsewhere exclude
this particular assignment. The coordinate-port obstruction in
\cref{s6:ports} gives a separate reason arbitrary degree-three LP circuits
cannot simply be replaced by degree-two pool-free gadgets.

Several restrictions do remove the missing aggregate information.
A fixed number of attachments permits the exact boundary construction
of \cref{s4:bounded-attachments}. Exact source supplies and a fixed
number of receiving products make total pool mass and attribute mass
affine boundary functions in \cref{s5:fixed-products}. With unbounded
feeds and outlets, exact contracts give the signed conservation system
of \cref{s5:quadratic-field}; fixed exceptions remain tractable even
with arbitrary attribute rank in \cref{s5:arbitrary-exceptions} when
common bounds are redundant. The support tests of
\cref{s5:restrictive-capacity} handle restrictive common throughput
bounds in the fully contracted scalar or affine-rank-one case, and
\cref{s5:throughput-optimization} also optimizes linear throughput
costs. Two full input-quality vectors admit a separate convex feasibility
reduction without a bypass degree restriction under the zero outlet
and common lower bounds of \cref{s5:two-vector-theorem}.
These results do not combine automatically: arbitrary exceptions together
with a restrictive common capacity, or unrestricted dense arc costs,
require an additional argument.

The unrestricted one-pool, one-quality bypass threshold question is
settled by \cref{s3:base-bypass,s3:upper-only,s3:constant-thm}; it is not
an open question here. Its tractable predecessors retain material
hypotheses: subdividing bypasses into artificial pools changes pool
count, and dropping source capacities changes shared resource coupling.
The surviving degree-two question fixes parameters absent from the
hardness constructions or imposes degree limits they do not meet.
Similarly, the two-quality response construction is an LP with interval
product specifications after penalties; it does not contradict the
exact-product-contract feasibility theorem.

The distinctions can now be stated concretely. Few nonlinear parameters
yield short basis-index certificates but do not alone yield polynomial
optimization. Independent bounded blocks, small retained boundaries,
and exact conservation identities each provide a different way to
control elimination. Exponentially many responses and strict local
optima obstruct particular descriptions or local arguments, while
\cref{s6:lp-hull} shows that they can coexist with exact polynomial
optimization. Any further classification must preserve the physical
couplings, objective class, contracts and arithmetic representation
that make these boundaries different.
